\chapter{Embedded Signal Processing on Arduino and Raspberry Pi}

\chapterline

Everything in this book so far has lived in a comfortable place: a notebook running NumPy, a
plot that redraws in a quarter second, arrays of \texttt{float64} that never overflow and never
run out. That is the right place to \emph{learn} DSP. It is not where DSP \emph{lives}. The
spectrogram that watches your forearm, the filter that smooths an accelerometer trace, the
peak-picker that decides a finger just tapped -- in the finished wrist controller those run on a
chip the size of a fingernail, powered by a coin cell, with no operating system, a few kilobytes
of RAM, and a hard deadline arriving every millisecond whether you are ready or not. The math is
identical. The engineering could not be more different.

This chapter is about that crossing -- taking the ideas we built on graph paper and making them
run, in real time, on a microcontroller (MCU) like an Arduino and on a small Linux computer like
a Raspberry Pi. We will see how to sample without jitter, why \texttt{float} sometimes becomes
your enemy and fixed-point your friend, how a ring buffer lets a finite chip process an infinite
stream, what it actually costs to run an FFT on a Cortex-M, and how processed numbers leave the
device as serial, BLE, or MIDI. The through-line goal -- a wrist-worn finger-gesture instrument --
forces every one of these choices to be concrete. There is a real-time budget, a real power budget,
and a real latency a musician will feel. By the end you should be able to look at the project and
say, with reasons, ``this part runs on the Arduino, that part runs on the Pi, and here is why.''

A word up front about the sensing, because it shapes every decision that follows. Finger movement
\emph{is} detectable from a wrist or forearm band, and we will treat it as a solved problem rather
than a wish. Two physical channels make it work. When you move a finger, the tendons in your
forearm slide and the muscle bellies change shape; a snug band of force sensors reads that
mechanical bulge -- this is \textbf{force myography} (FMG). At the same time the muscles fire
electrical pulses that a pair of skin electrodes can pick up -- \textbf{surface electromyography}
(sEMG). Neither one needs a camera, a glove, or a machine-learning model. The method is direct: the
user runs a short guided \emph{enrollment}, performing each finger movement a few times while the
phone records the signal signature; from those recordings we build a per-user template for each
finger. Live detection is then a match -- compare the current signal to the calibrated templates
and fire the closest one. Calibration, not classification, is the heart of the instrument, and the
whole chapter builds toward making that loop run on a chip you can wear.

\section{The shape of the problem: streaming, not batch}
\label{ch40:streaming}

On a laptop you think in \emph{batches}. You have the whole signal in an array, you call
\texttt{np.fft.fft(x)}, you get the answer. On an embedded device the signal does not exist yet.
It is being born one sample at a time, forever, and you must produce outputs while later inputs
are still arriving. The mental model is a pipe, not a bucket.

\begin{intuitionbox}
Batch DSP is a photograph: the whole signal is laid out, and you operate on it at leisure.
Streaming DSP is a river you stand in: water arrives at a fixed rate, you must scoop and act
before it flows past, and you never get to hold all of it at once. Every embedded design decision
-- buffer sizes, fixed-point, interrupt timing -- exists to keep up with the river without falling
behind and without drowning in stored water.
\end{intuitionbox}

This single shift reshapes everything. You cannot store an infinite stream, so you keep a short
sliding window (a \emph{ring buffer}). You cannot afford a sloppy sample clock, because a filter
designed for $f_s = 1000$~Hz quietly lies if the true rate wobbles. You cannot pause to garbage-
collect or wait on a disk. And you have a \textbf{real-time budget}: if your per-sample (or
per-block) work takes longer than the time between samples (or blocks), the river overflows and
you lose data. That budget is the spine of the whole chapter.

\begin{focusbox}
Three numbers govern every embedded DSP design, and you should be able to state all three for any
project before writing code:
\begin{enumerate}
  \item \textbf{Sample rate} $f_s$ -- how often a new sample arrives (the deadline clock).
  \item \textbf{Work per sample/block} -- how many cycles your processing costs.
  \item \textbf{Cycle budget} -- the MCU clock times the time between deadlines.
\end{enumerate}
If work per block exceeds the cycle budget, no amount of clever code saves you; you must lower
$f_s$, shrink the algorithm, or move to a faster processor.
\end{focusbox}

\section{Getting samples in: timers, ADCs, and the enemy called jitter}
\label{ch40:sampling}

Sampling theory told us \emph{how often} to sample. Embedded reality adds a brutal footnote: the
samples must land on an \emph{evenly spaced} time grid. The whole edifice of the DFT, every filter
coefficient, every frequency you read off a spectrum -- all of it assumes that sample $n$ happened
at exactly $t_n = n/f_s$. When the actual instants drift around that grid, we call the error
\textbf{jitter}, and it is one of the most insidious problems in embedded DSP because the code
looks correct and the numbers are merely wrong.

\subsection{Where samples come from}

A microcontroller has a built-in ADC. The naive way to use it is a loop:

\begin{lstlisting}[language=C, caption={The WRONG way to sample: timing set by loop speed.}]
void loop() {
  int x = analogRead(A0);   // read sensor
  process(x);               // do DSP
  delay(1);                 // "1 ms" -- but is it?
}
\end{lstlisting}

This is a trap. The interval between samples is now \texttt{delay(1)} \emph{plus} however long
\texttt{analogRead} took \emph{plus} however long \texttt{process} took \emph{plus} whatever
interrupt fired in between. The sample spacing breathes with your code. If \texttt{process}
sometimes branches into a slower path, your $f_s$ changes mid-stream.

\begin{trapbox}
\texttt{delay()}-based sampling produces a sample rate that depends on your processing time. The
moment your DSP load varies (a conditional, a heavier block every $N$ samples, an occasional
log-print), the time grid warps. Spectra smear, filters detune, and the bug is nearly invisible
because every individual sample looks fine. \textbf{Never set your sample clock from inside the
work loop.}
\end{trapbox}

The right way is to let a hardware \textbf{timer} fire an interrupt at a fixed rate and have the
interrupt -- and \emph{only} the interrupt -- take the sample. The timer counts the crystal
oscillator directly; it does not know or care what your code is doing. This decouples the sample
clock from the workload.

\begin{lstlisting}[language=C, caption={Timer-driven sampling on an AVR Arduino (Timer1, $f_s=1$~kHz).}]
volatile int16_t latest;
volatile bool ready;

ISR(TIMER1_COMPA_vect) {       // fires exactly every 1 ms
  latest = analogRead(A0);     // sample on the grid
  ready = true;                // flag the main loop
}

void setup() {
  cli();
  TCCR1A = 0; TCCR1B = 0; TCNT1 = 0;
  OCR1A  = 249;                // 16 MHz / 64 / (249+1) = 1000 Hz
  TCCR1B |= (1 << WGM12);      // CTC mode: clear timer on compare
  TCCR1B |= (1 << CS11) | (1 << CS10);  // prescaler = 64
  TIMSK1 |= (1 << OCIE1A);     // enable compare interrupt
  sei();
}

void loop() {
  if (ready) { ready = false; process(latest); }
}
\end{lstlisting}

Now the sample instants are nailed to the crystal. The main loop can be late to \emph{process} a
sample, but the sample itself was taken on time. (Better still on richer chips: have the timer
trigger the ADC in hardware and have the ADC's done-interrupt fetch the result, so even the
\texttt{analogRead} latency leaves the timing path. On a Cortex-M you would wire a timer to the
ADC trigger and the ADC to DMA, and the CPU only wakes when a whole block is ready.)

\subsection{What jitter does to a spectrum}

Jitter is not just untidy -- it injects a noise floor. Suppose each true instant $t_n = n/f_s$ is
perturbed to $t_n + \delta_n$ with $\delta_n$ small and random. For a sinusoid
$x(t)=A\sin(2\pi f t)$ the sampled value is
\[
  x(t_n+\delta_n) \approx A\sin(2\pi f t_n) + \underbrace{2\pi f A\cos(2\pi f t_n)}_{\text{slope}}\,\delta_n .
\]
The error term scales with $f$: \emph{the higher the frequency, the more a fixed timing error
costs you}. A timing wobble of $\delta$ seconds turns into an amplitude error proportional to
$2\pi f \delta$. This is why high-frequency, high-resolution sampling is so demanding on the
clock, and why jitter shows up as a raised, roughly white noise floor that worsens toward the top
of the band.

\begin{figure}[H]
\centering
\begin{tikzpicture}
\begin{axis}[
  width=12cm, height=5cm,
  xlabel={frequency (Hz)}, ylabel={magnitude (dB)},
  xmin=0, xmax=200, ymin=-70, ymax=5,
  axis lines=left, enlargelimits=false,
  legend pos=north east, legend cell align=left,
]
\addplot[sigblue, thick] coordinates {(0,-65)(40,-65)(49,-60)(50,0)(51,-60)(60,-66)(200,-66)};
\addlegendentry{clean clock}
\addplot[specred, thick, dashed] coordinates {(0,-50)(40,-46)(49,-40)(50,0)(51,-40)(70,-42)(120,-38)(200,-32)};
\addlegendentry{jittered clock}
\end{axis}
\end{tikzpicture}
\caption{Sampling a single 50~Hz tone. A clean timer keeps the noise floor low and flat (solid).
Clock jitter (dashed) raises the floor, and the rise grows toward higher frequencies because a
fixed timing error perturbs fast-moving samples more.}
\label{ch40:fig-jitter}
\end{figure}

\begin{whybox}
A hardware timer fires from a counter clocked by the crystal oscillator, completely independent
of what the CPU is executing. Its edges are as regular as the crystal itself (tens of ppm). A
software \texttt{delay()} loop, by contrast, measures time in instructions executed, so anything
that changes the instruction count -- a branch, an interrupt, a cache miss on bigger chips --
moves the sample. Timers work because they move the timing decision out of software entirely and
into dedicated silicon.
\end{whybox}

\begin{goalbox}
The wrist controller streams several channels at once: a MEMS IMU (accelerometer and gyro), an FMG
array of force sensors reading tendon bulge, and one or two sEMG channels reading muscle pulses. For
gesture timing a musician can feel, you want low, \emph{stable} latency, so every sample clock must
be rock-steady. Practically: configure the IMU's own internal sample-rate generator (a good MEMS
IMU has one, e.g. an ICM-20948 at a programmed output data rate) and read it over its FIFO, rather
than polling \texttt{analogRead} in \texttt{loop()}. The IMU's clock becomes your time grid. For the
analog FMG and EMG channels, drive the MCU ADC from a timer as above so they land on the same grid
and can be compared sample for sample. Either way, the rule holds: the clock comes from hardware,
never from the work loop. Sample-aligned channels are what let the calibration step later treat one
finger movement as a single multi-sensor signature.
\end{goalbox}

\section{Ring buffers: a finite window onto an infinite stream}
\label{ch40:ring}

You cannot store the whole stream, but most DSP only needs a short recent history: an FIR filter
of length $M$ needs the last $M$ samples; a 256-point FFT needs the last 256; an envelope detector
needs maybe a few dozen. The data structure that holds ``the last $N$ samples'' efficiently is the
\textbf{ring buffer} (circular buffer): a fixed array with a write index that wraps around to the
start when it hits the end. Geometrically it is a loop of tape -- you keep writing forward, and the
oldest sample is overwritten by the newest, so the buffer always holds a window sliding along the
signal.

\begin{figure}[H]
\centering
\begin{tikzpicture}[scale=1]
\def\Rad{1.6}
\foreach \i/\a in {0/90,1/45,2/0,3/-45,4/-90,5/-135,6/180,7/135} {
  \node[draw, circle, minimum size=7mm, fill=fillsig!40] (s\i) at (\a:\Rad) {\footnotesize \i};
}
\draw[->, thick, structurepurple] (135:\Rad+0.5) arc (135:95:\Rad+0.5);
\node[structurepurple] at (90:\Rad+1.05) {\footnotesize write here, then wrap};
\node at (0,0) {\footnotesize 8-slot ring};
\end{tikzpicture}
\caption{A ring buffer of 8 slots. The write index advances around the loop; when it passes slot 7
it returns to slot 0, overwriting the oldest sample. The array index is the sample number taken
modulo the buffer length.}
\label{ch40:fig-ring}
\end{figure}

The key arithmetic is that the physical array index is the logical sample number \emph{modulo} the
buffer length $N$. If $N$ is a power of two, that modulo is a single bitwise AND -- one cycle --
which is exactly why DSP buffer sizes are almost always powers of two.

\begin{lstlisting}[language=C, caption={A power-of-two ring buffer with a single-cycle wrap.}]
#define N 256                 // must be a power of two
int16_t buf[N];
uint16_t head = 0;            // next write position

static inline void ring_push(int16_t x) {
  buf[head] = x;
  head = (head + 1) & (N - 1);   // wrap via AND, no modulo divide
}

// read the sample k steps back (k=0 is newest)
static inline int16_t ring_get(uint16_t k) {
  return buf[(head - 1 - k) & (N - 1)];
}
\end{lstlisting}

\begin{happeningbox}
A ring buffer turns ``the last $N$ samples'' into a fixed footprint that never grows. No
allocation, no shifting of array elements, no copying -- pushing a sample is one store and one
masked increment. The illusion of a smoothly sliding window is created by moving the \emph{index}
instead of moving the \emph{data}. This is the single most important data structure in real-time
DSP, and it is why an MCU with 4~KB of RAM can process a signal that has been running for hours.
\end{happeningbox}

\begin{trapbox}
A ring buffer shared between an interrupt (writer) and the main loop (reader) is a classic
concurrency hazard. The simplest safe pattern: the ISR only \emph{pushes}, the main loop only
\emph{reads}, indices are a single aligned word (atomic to read/write on the MCU), and shared
flags are \texttt{volatile}. For block processing, use \emph{two} buffers (double buffering / ping-
pong): the ISR fills buffer~A while the main loop processes buffer~B, then they swap. Never let the
reader and writer touch the same buffer at the same time.
\end{trapbox}

\section{Numbers on a small chip: fixed-point vs floating-point}
\label{ch40:fixed-point}

On your laptop every number is a 64-bit float and you never think about it. On a small MCU you
must, because the way you represent numbers directly sets your speed, your memory, and your
correctness. There are two camps: floating-point and fixed-point.

\subsection{What floating-point costs}

A floating-point number stores a sign, an exponent, and a mantissa -- effectively scientific
notation in binary, $\pm m \times 2^{e}$. Its glory is \emph{dynamic range}: it represents both
$0.0001$ and $10000$ with similar relative precision, so you rarely overflow and rarely lose all
your bits. The catch on embedded hardware is the \textbf{FPU} (floating-point unit). A chip with a
hardware FPU (Cortex-M4F, M7, the Raspberry Pi) does a float multiply in a cycle or few. A chip
\emph{without} one (the classic 8-bit AVR Arduino Uno, Cortex-M0/M0+) must emulate every float
operation in software -- dozens to hundreds of cycles each. On an Uno, a single \texttt{float}
multiply can cost what an integer multiply costs many times over.

\subsection{What fixed-point is}

Fixed-point keeps the speed of integer arithmetic but reintroduces fractions by \emph{agreeing on
where the binary point sits}. In the common \textbf{Q-format} notation, $Q_{m.n}$ means an integer
interpreted as having $n$ fractional bits: the stored integer $X$ represents the real value
$X / 2^{n}$. For example, in $Q_{1.15}$ (often just ``Q15'') a 16-bit signed integer represents a
real number in $[-1, 1)$ with $2^{15}$ steps; the value $0.5$ is stored as $16384$.

\begin{intuitionbox}
Floating-point lets the binary point \emph{float} to wherever the number needs it -- great range,
hardware cost if there is no FPU. Fixed-point nails the binary point to one agreed place -- you
trade away automatic range for the raw speed of plain integers. Choosing between them is choosing
\emph{who tracks the scale}: the silicon (float) or you, the programmer (fixed).
\end{intuitionbox}

Arithmetic in Q-format follows simple rules once you track the point:
\begin{itemize}
  \item \textbf{Add/subtract:} operands must share the same $n$; the result keeps that $n$. Watch
    for overflow -- two Q15 numbers near $1$ sum to near $2$, which no longer fits.
  \item \textbf{Multiply:} $Q_{a}\times Q_{b}$ gives a product with $a+b$ fractional bits. Two Q15
    values multiplied give a 32-bit Q30 result; to return to Q15 you shift right by 15. You must
    multiply into a \emph{wider} accumulator (16$\times$16 $\to$ 32) or you lose the top bits.
\end{itemize}

\begin{lstlisting}[language=C, caption={A Q15 fixed-point multiply with a 32-bit accumulator.}]
typedef int16_t q15_t;        // value v stored as round(v * 32768)

static inline q15_t q15_mul(q15_t a, q15_t b) {
  int32_t p = (int32_t)a * (int32_t)b;   // Q30 in 32 bits
  return (q15_t)(p >> 15);               // back to Q15 (truncating)
}
// example: 0.5 * 0.5 -> 16384 * 16384 = 268435456; >>15 = 8192 = 0.25
\end{lstlisting}

\begin{trapbox}
Two fixed-point sins dominate embedded bug reports. \textbf{Overflow}: a Q15 accumulator quietly
wraps from $+0.99$ to $-1.0$ when a filter's partial sums exceed the range; the fix is a wider
accumulator (Q31) and \emph{saturating} arithmetic that clamps instead of wrapping. \textbf{Loss
of precision}: shifting right too early throws away low bits; keep full width through the whole
inner sum and round/shift only at the very end. CMSIS-DSP's functions are written precisely to
manage these for you, which is a strong reason to use them rather than rolling your own.
\end{trapbox}

\begin{connectionbox}
Quantization and fixed-point are two faces of the same coin from the ADC chapter. Quantization
error is the gap between the true value and the nearest representable rung; in $Q_{m.n}$ that rung
spacing is exactly $2^{-n}$. The ``6~dB per bit'' rule is the same here: each extra fractional bit
halves the step and buys about 6~dB of headroom against round-off noise. The noise you analyzed
statistically in the noise chapter is the same round-off noise your fixed-point filter generates
internally -- which is why understanding random signals pays off the moment you go embedded.
\end{connectionbox}

\subsection{A practical rule of thumb}

If your chip has an FPU, write in \texttt{float} (or even \texttt{double} on a Pi) first. Floats
are easier to get correct, and on an M4F/M7 they are fast. Reach for fixed-point only when you
have no FPU (AVR, M0), when you need every last cycle or microamp, or when a vectorized integer
path (CMSIS-DSP Q15) is meaningfully faster than the float path on your part. Premature fixed-point
is a top source of subtle, hard-to-find DSP bugs.

\section{Running real DSP: filters and FFTs on an MCU vs a Pi}
\label{ch40:running}

Now the heart of it. We have samples on a clean grid, stored in a ring buffer, in a numeric format
the chip likes. What does it cost to actually \emph{do} the DSP?

\subsection{FIR and IIR filters on the MCU}

A direct-form FIR filter of length $M$ is $y[n]=\sum_{k=0}^{M-1} h[k]\,x[n-k]$ -- $M$ multiply-
accumulates (MACs) per output sample. This maps beautifully onto a ring buffer: the taps multiply
the last $M$ ring entries. A Cortex-M4 has a single-cycle MAC, so an $M=64$ FIR costs roughly $64$
cycles per sample; at $f_s=1$~kHz that is $64{,}000$ cycles/s, a rounding error against an
$80$~MHz clock. The same filter on an 8-bit AVR in floating point could cost a hundred times more
per tap and might not keep up.

An IIR biquad is even cheaper -- five MACs per sample per second-order section -- and is the
workhorse for the gentle low-pass and band-pass filters you will run on IMU and EMG signals. The
price of that efficiency is sensitivity to coefficient quantization and the need to watch internal
overflow, which is exactly why biquads are usually run as cascaded second-order sections in
fixed-point.

\begin{lstlisting}[language=C, caption={A single biquad (Direct Form II transposed), float version.}]
typedef struct { float b0,b1,b2,a1,a2,z1,z2; } biquad_t;

static inline float biquad(biquad_t* s, float x) {
  float y = s->b0 * x + s->z1;
  s->z1   = s->b1 * x - s->a1 * y + s->z2;
  s->z2   = s->b2 * x - s->a2 * y;
  return y;                         // 5 mul, 4 add, 2 state updates
}
\end{lstlisting}

\subsection{The FFT: CMSIS-DSP on the MCU}

For an FFT on an ARM Cortex-M, do not write your own -- use \textbf{CMSIS-DSP}, ARM's free,
hand-optimized DSP library. It provides radix-based FFTs in float (\texttt{arm\_rfft\_fast\_f32})
and fixed-point (\texttt{arm\_rfft\_q15}), plus FIR/IIR, statistics, and vector math, all tuned to
use the chip's MAC and (on M4/M7) SIMD instructions.

\begin{lstlisting}[language=C, caption={A 256-point real FFT and magnitude with CMSIS-DSP.}]
#include "arm_math.h"
#define NFFT 256
arm_rfft_fast_instance_f32 fft;
float in[NFFT], spec[NFFT], mag[NFFT/2];

void fft_setup(void) { arm_rfft_fast_init_f32(&fft, NFFT); }

void fft_block(void) {
  arm_rfft_fast_f32(&fft, in, spec, 0);    // forward real FFT
  arm_cmplx_mag_f32(spec, mag, NFFT/2);    // |X[k]| for each bin
}
\end{lstlisting}

How long does that take? Order-of-magnitude: an $N=256$ float FFT on an $80$~MHz Cortex-M4F is on
the order of tens of microseconds. If you want a new spectrum every $10$~ms (a $100$~Hz frame rate,
plenty for gesture work), the FFT eats well under $1\%$ of the budget. This is the headline result
of the chapter: \emph{modern MCUs are genuinely fast enough to do real frequency-domain DSP in
real time}, as long as you respect the budget arithmetic.

\begin{walkbox}
A 256-point FFT in $\sim 30~\mu$s. A frame every $10$~ms ($10{,}000~\mu$s). What fraction of the
processor does the FFT use, and how many such FFTs could you run per frame before you run out of
time? \\[2pt]
\emph{Answer:} $30/10000 = 0.3\%$ per FFT. The math alone would allow on the order of a few hundred
such FFTs per frame -- but in practice you would hit memory and I/O limits, and windowing/copying
overhead, long before the FFT arithmetic itself becomes the bottleneck.
\end{walkbox}

\subsection{The Pi side: NumPy and SciPy}

A Raspberry Pi is a different animal -- a full Linux computer with a multi-core ARM CPU, hundreds
of megabytes to gigabytes of RAM, and an FPU on every core. Here you write Python with NumPy and
SciPy exactly as you did while learning. \texttt{np.fft.rfft}, \texttt{scipy.signal.sosfilt},
\texttt{scipy.signal.spectrogram} -- all available, all fast (NumPy's FFT and filters are C/Fortran
under the hood). The cost you trade for that comfort is real-time \emph{determinism}: Linux is not
a hard real-time OS. Most of the time your block runs in a millisecond; occasionally the scheduler,
a page fault, or another process makes it late. For a feel-it-instantly instrument that occasional
hiccup matters.

\begin{lstlisting}[language=Python, caption={Streaming block processing on a Raspberry Pi.}]
import numpy as np
from scipy.signal import butter, sosfilt

fs = 1000
sos = butter(4, 20, 'low', fs=fs, output='sos')  # 20 Hz low-pass
zi  = np.zeros((sos.shape[0], 2))                 # filter state

def process_block(block):           # block: 256 new samples
    global zi
    y, zi = sosfilt(sos, block, zi=zi)            # stateful -> streaming
    spec  = np.abs(np.fft.rfft(y * np.hanning(len(y))))
    return y, spec
\end{lstlisting}

\begin{aestheticbox}
The same five-line biquad is, on the MCU, a hand-managed struct of floats you must budget cycles
for, and on the Pi, a \texttt{butter()} call you barely think about. Neither is ``better.'' The
MCU gives you determinism and microamps at the cost of effort; the Pi gives you effortless
horsepower at the cost of determinism and milliwatts. Mature embedded design is knowing which
currency you are spending and choosing on purpose.
\end{aestheticbox}

\section{Getting answers out: serial, BLE, and MIDI}
\label{ch40:output}

Processing is half the job; the result has to leave the device. Three channels matter for this
project.

\textbf{Serial (UART/USB).} The workhorse for development and wired use. Simple, reliable, and on
USB fast enough for raw sample streams. You will use it constantly for debugging -- but printing
inside your DSP loop is itself a timing hazard (see the trap below).

\textbf{BLE (Bluetooth Low Energy).} The right wireless link for a wearable: low power, ubiquitous
on phones and laptops, and -- crucially -- it supports \textbf{BLE-MIDI}, a standard profile that
carries MIDI messages over Bluetooth. A board like a Nordic nRF52 or an ESP32 can advertise as a
BLE-MIDI device that a DAW or synth app sees as a normal MIDI controller, with no dongle.

\textbf{MIDI.} The destination format for our instrument. MIDI is delightfully thin: most messages
are three bytes -- a status byte (command + channel) and two data bytes. A Note-On is
\texttt{0x90 | channel}, then note number ($0$--$127$), then velocity ($0$--$127$). A continuous
controller (CC), perfect for mapping a gesture to a parameter like filter cutoff or vibrato, is
\texttt{0xB0 | channel}, controller number, value.

\begin{lstlisting}[language=C, caption={Emitting MIDI: a note and a continuous-control message.}]
void midi_note_on(uint8_t ch, uint8_t note, uint8_t vel) {
  uint8_t m[3] = { (uint8_t)(0x90 | (ch & 0x0F)), note, vel };
  midi_send(m, 3);            // over USB-MIDI, serial-MIDI, or BLE-MIDI
}
void midi_cc(uint8_t ch, uint8_t cc, uint8_t val) {
  uint8_t m[3] = { (uint8_t)(0xB0 | (ch & 0x0F)), cc, val };
  midi_send(m, 3);
}
\end{lstlisting}

\begin{trapbox}
\texttt{Serial.print()} for debugging inside a sensor- or audio-rate loop can block for
milliseconds while the UART drains its buffer, blowing your real-time deadline and corrupting your
sample timing. For diagnostics, either print from a \emph{lower-priority} context (the slow main
loop, not the ISR), throttle to a few messages per second, toggle a GPIO pin and watch it on a
logic analyzer, or write to a ring buffer and flush it when idle. Logging is I/O, and I/O is the
enemy of determinism.
\end{trapbox}

\begin{goalbox}
The wrist controller's output path is the whole point: finger gestures in, music out. The natural
design is BLE-MIDI -- the wearable advertises as a standard MIDI controller, and any phone, tablet,
or laptop running a synth receives it with zero custom software. Map \emph{discrete} finger
movements (a per-finger tap or press, matched to its calibrated template) to Note-On/Note-Off, and
\emph{continuous} signals (finger pressure from FMG, wrist tilt) to CC messages so the player can
bend, swell, and modulate. Keep total latency -- sample to MIDI out -- as low as you can; beyond
about $20$--$30$~ms a player begins to feel the lag, which is why the capstone leans on a fast
onset detector (Section~\ref{ch40:worked}).
\end{goalbox}

\section{Power and real-time budgets}
\label{ch40:budgets}

A wearable lives on a battery, so two budgets run in parallel: cycles (can I finish in time?) and
joules (how long until it dies?). They pull in opposite directions, and good design balances them.

\subsection{The real-time budget, made concrete}

The rule is unforgiving: \emph{work per deadline must be less than time per deadline}. Suppose you
sample an IMU at $f_s = 200$~Hz and run the heavier feature/analysis block on $32$ samples, so a
block arrives every $32/200 = 160$~ms. (This block sets the analysis cadence, not the play latency
-- the fast onset path of Section~\ref{ch40:worked} fires far sooner.) On an $80$~MHz MCU that is
$80{,}000{,}000 \times 0.160 = 12.8$~million cycles of budget per block -- enormous. Your filtering and a small FFT might use a few tens of thousands.
You are nowhere near the cliff, which is typical for motion-rate DSP and is exactly why an MCU
suffices. The budget only gets tight at audio rates ($f_s$ in the tens of kHz) or with big FFTs.

\begin{figure}[H]
\centering
\begin{tikzpicture}
\begin{axis}[
  width=12cm, height=5cm,
  xlabel={per-block work (millions of cycles)}, ylabel={fraction of budget},
  xmin=0, xmax=14, ymin=0, ymax=1.3,
  xtick={0,2,4,6,8,10,12,14},
  ytick={0,0.5,1.0}, axis lines=left, enlargelimits=false,
]
\addplot[sigblue, thick, domain=0:12.8, samples=2] {x/12.8};
\addplot[specred, dashed, domain=0:14, samples=2] {1};
\node[specred, anchor=south east] at (axis cs:14,1.0) {\footnotesize deadline (100\%)};
\addplot[targetgreen, only marks, mark=*] coordinates {(0.05,0.004)};
\node[targetgreen, anchor=west] at (axis cs:0.4,0.12) {\footnotesize typical motion-DSP load};
\end{axis}
\end{tikzpicture}
\caption{Real-time budget for a $200$~Hz IMU processed in $32$-sample blocks on an $80$~MHz MCU:
$12.8$M cycles available per block. As work rises toward the dashed deadline you approach the
cliff; motion-rate DSP (green dot) sits comfortably near the origin.}
\label{ch40:fig-budget}
\end{figure}

\subsection{The power budget}

Battery life is set by average current. The dominant lever on an MCU is \emph{sleep}: a Cortex-M
draws tens of milliamps running flat out but only microamps in a deep-sleep mode. Because motion
DSP uses a tiny fraction of the cycle budget, the chip should compute its block, then \emph{sleep}
until the next timer or IMU interrupt. This ``compute-and-sleep'' duty cycling can cut average
current by one to two orders of magnitude.

\begin{lstlisting}[language=C, caption={Compute-then-sleep duty cycling (pseudocode).}]
for (;;) {
  if (block_ready) {          // set by the IMU FIFO interrupt
    block_ready = false;
    process_block();          // a few tens of microseconds of work
    maybe_emit_midi();
  }
  enter_low_power_sleep();    // wake on next interrupt
}
\end{lstlisting}

\begin{applicationbox}
The power gap between Arduino-class and Pi-class hardware is stark. A sleeping nRF52 doing
intermittent motion DSP and BLE can run for \emph{days} on a small LiPo. A Raspberry Pi (even a
Zero) idles in the hundreds of milliamps and has no microamp sleep -- it is a poor wearable, often
needing recharging within hours and a far bulkier battery. This single fact, more than raw compute,
is why a battery-powered wrist device gravitates to an MCU and a tethered or desktop processor
gravitates to a Pi.
\end{applicationbox}

\begin{connectionbox}
The power story loops back to sampling theory. Lower $f_s$ means fewer interrupts, fewer wake-ups,
and less radio traffic -- so the Nyquist choice is also a battery choice. Sample only as fast as
the signal demands: human gesture bandwidth is a few tens of hertz, so a $100$--$200$~Hz IMU rate
is plenty, and pushing higher just burns joules for no extra information.
\end{connectionbox}

\section{Arduino vs Raspberry Pi for the wrist controller}
\label{ch40:compare}

Time to decide. The two platforms are not competitors so much as different tools, and the right
answer for this project is usually ``use both'' -- a thin wrist node and a capable phone.

\begin{center}
\small
\begin{tabular}{@{}p{2.6cm}p{5.2cm}p{5.2cm}@{}}
\toprule
 & \textbf{Arduino / MCU} (nRF52, ESP32, M4) & \textbf{Raspberry Pi} (Linux SBC) \\
\midrule
Compute & MHz, single core, KB--MB RAM & GHz, multi-core, GB RAM \\
Numerics & fixed-point or FPU float; CMSIS-DSP & full \texttt{float}/\texttt{double}; NumPy/SciPy \\
Real time & hard, deterministic (ISR + timer) & soft; OS jitter, no hard deadlines \\
Power & mA active, $\mu$A asleep -- days on a coin/LiPo & hundreds of mA -- hours, bulky battery \\
I/O & native ADC, SPI/I2C, BLE-MIDI & USB, Wi-Fi/BT, needs ADC HAT for analog \\
Dev effort & C, manual budgets, flash-and-test & Python, fast iteration, easy logging \\
Best at & always-on sensing, wireless, low latency & calibration, conditioning, deep offline analysis \\
\bottomrule
\end{tabular}
\end{center}

\begin{buildbox}
\textbf{Recommended split for the wrist controller.} The architecture is deliberately lopsided: a
\emph{thin} wrist node and a \emph{capable} phone.
\begin{itemize}
  \item \textbf{On the wrist (MCU, e.g. nRF52840):} read the IMU FIFO and the timer-driven FMG/EMG
    channels on hardware clocks, apply only light conditioning -- a couple of biquads and an
    envelope -- compress, and BLE-stream the channels. For live play, run a lightweight on-device
    template match (a short correlation against the few calibrated per-finger templates) so the
    note fires with minimal latency even if the phone is busy. Compute-and-sleep between blocks for
    battery life. This is the always-on, low-latency, low-power core.
  \item \textbf{On the phone (or a Pi/laptop):} the heavy lifting lives here. The phone runs the
    guided calibration routine, builds and stores the per-finger templates and thresholds,
    conditions the streamed signals, and does deep offline analysis from the logged streams to
    refine the templates. There is no learned model -- ``analysis'' here means filtering, alignment,
    averaging the enrollment repetitions into a clean template, and picking thresholds. The phone
    then ships the finished templates back down to the wrist node for the live path.
\end{itemize}
The general principle: keep the wrist node cheap in cycles and joules, and let the phone -- which
has power to spare and a screen to guide the user -- own calibration and conditioning. The wrist
sleeps; the phone thinks.
\end{buildbox}

\begin{goalbox}
\textbf{What the wrist can and cannot sense.} Be precise here, because the boundary is the whole
design. \emph{Which finger just moved} is observable from the forearm: tendon bulge (FMG) and
muscle pulses (sEMG) carry a distinct per-finger signature, and after a short per-user calibration a
template match reads it reliably. WristFlex demonstrated FMG-based finger-tap discrimination from a
band of force sensors; Thalmic Labs' Myo and the open MyoWare front end read the same intent from
sEMG. So per-finger \emph{detection} -- this finger tapped, that finger pressed, this one is held
down -- is firmly on the doable side, and it is plenty for playing an instrument.

The genuinely hard part, and the one limit worth stating plainly, is continuous per-\emph{joint}
finger \emph{angles}. You cannot recover them from a wrist or forearm band: the mapping from finger pose to
wrist-surface signal is underdetermined, and the bulk motion of the wrist simply does not contain
that information. Recovering smooth joint angles needs sensors on the fingers themselves --
instrumented rings, as in the finger-worn ``Learning on the Rings'' tracking work, or a glove. Treat
that as an upgrade path, not a launch requirement. Everything the project actually needs -- detect
which finger moved, how hard, and when -- sits comfortably on the achievable side.
\end{goalbox}

\begin{historybox}
The split we just drew is an old one in disguise. In the 1980s the Texas Instruments TMS320 fixed-
point DSP chips earned their living doing exactly this -- a small, deterministic processor crunching
filters and FFTs in real time on a tight power budget -- while general-purpose CPUs did the
flexible, non-real-time work. ARM's CMSIS-DSP library and the Cortex-M's single-cycle MAC are the
direct descendants of that fixed-point DSP lineage, now small and cheap enough to wear on your
wrist. The Raspberry Pi, meanwhile, is the heir to the workstation. Embedded DSP has always been
the art of putting the right work on the right side of that line.
\end{historybox}

\section{The capstone: a buildable finger-gesture controller}
\label{ch40:worked}

Now assemble everything into one concrete, end-to-end project (Figure~\ref{ch40:fig-arch}). The
instrument detects which finger moved, how hard, and when, and emits MIDI -- no camera, no machine
learning, just multimodal sensing plus per-finger calibration. This section is meant to be
\emph{built}.

\begin{figure}[H]
\centering
\begin{tikzpicture}[
  node distance=6mm,
  box/.style={draw, rounded corners, align=center, font=\footnotesize, minimum height=8mm, inner sep=3pt},
  >=Stealth,
]
\node[box, fill=fillsig!25] (sensors) {IMU + FMG array\\+ sEMG + fast accel};
\node[box, fill=fillsig!25, right=of sensors] (node) {wrist node (MCU)\\condition, match,\\compress, sleep};
\node[box, fill=green!12, right=14mm of node] (phone) {phone\\calibration, conditioning,\\offline analysis};
\node[box, fill=orange!15, below=10mm of phone] (synth) {synth / DAW};
\draw[->] (sensors) -- (node);
\draw[->] ([yshift=2mm]node.east) -- node[above, font=\scriptsize]{BLE stream} ([yshift=2mm]phone.west);
\draw[<-] ([yshift=-2mm]node.east) -- node[below, font=\scriptsize]{templates} ([yshift=-2mm]phone.west);
\draw[->] (node) |- node[pos=0.75, above, font=\scriptsize]{BLE-MIDI (live)} (synth.west);
\end{tikzpicture}
\caption{System architecture. A thin wrist node conditions and template-matches the fused sensor
channels and emits BLE-MIDI for the low-latency live path. It BLE-streams raw channels to the phone,
which runs calibration and deep offline analysis and ships finished per-finger templates back down.}
\label{ch40:fig-arch}
\end{figure}

\subsection{Bill of materials and wiring}

\begin{itemize}
  \item \textbf{MCU:} Nordic nRF52840 (Cortex-M4F, FPU, BLE 5, deep-sleep in microamps). The wrist
    node.
  \item \textbf{IMU:} ICM-20948 (or LSM6DSx) over SPI, programmed to a $200$~Hz output data rate
    with FIFO and a data-ready interrupt line.
  \item \textbf{High-rate accelerometer:} a second, fast accelerometer (a few kHz) used as a
    \emph{contact / bio-acoustic} sensor -- it hears the tendon snap and skin tap that a slow IMU
    misses, the trick ViBand (UIST~'16) used to turn a watch into a bio-acoustic input device.
  \item \textbf{FMG array:} four to eight force-sensitive resistors (FSRs) in a band pressed against
    the underside of the forearm, each into an ADC channel through a voltage divider.
  \item \textbf{sEMG:} one or two MyoWare (or equivalent) sEMG front-end boards, each a single
    analog channel into the ADC.
  \item \textbf{Power:} a small LiPo (e.g. $150$~mAh) and a charge/regulator board.
\end{itemize}

Wiring is straightforward: IMU on SPI with its interrupt to a GPIO; the high-rate accelerometer on a
second SPI or I\textsuperscript{2}C bus with its own interrupt; the FSR dividers and sEMG outputs to
ADC pins, all sampled from one hardware timer so they share a time grid.

\subsection{The pipeline, channel by channel}

\begin{enumerate}
  \item \textbf{Acquire (hardware clocks).} The IMU FIFO interrupt and the ADC timer drive
    acquisition; the high-rate accelerometer interrupts separately. No \texttt{delay()}, no polling.
    Each channel lands on its own steady grid.
  \item \textbf{Condition.} Run each slow channel (IMU, FMG) through a biquad low-pass near $20$~Hz
    to kill motion noise; band-pass the sEMG ($\sim 20$--$450$~Hz) and take its envelope (rectify
    then low-pass), since the useful sEMG signal is the \emph{amplitude} of the muscle burst, not its
    waveform. Five MACs per biquad per sample.
  \item \textbf{Feature.} Per short window, form a small feature vector: FMG channel levels, sEMG
    envelope levels, a high-rate-accelerometer onset/energy spike, and a couple of IMU stats. This
    fixed-length vector is the live ``signature.''
  \item \textbf{Match.} Compare the live signature against the calibrated per-finger templates
    (next subsection) -- nearest template by correlation or distance, gated by an onset threshold so
    silence stays silent.
  \item \textbf{Emit.} Map the matched finger and its pressure to MIDI -- Note-On for a tap, a held
    note while pressed, CC for continuous pressure or wrist tilt -- and send over BLE-MIDI.
  \item \textbf{Sleep.} Enter low-power mode until the next interrupt.
\end{enumerate}

\begin{figure}[H]
\centering
\begin{tikzpicture}
\begin{axis}[
  width=12cm, height=6.2cm,
  xlabel={time (ms)}, ylabel={channel (offset)},
  xmin=0, xmax=200, ymin=-0.5, ymax=6.5,
  ytick={1,3,5}, yticklabels={accel, sEMG env, FMG},
  axis lines=left, enlargelimits=false, samples=120, domain=0:200,
  legend pos=north east, legend cell align=left,
]
% FMG: slow rise and fall (tendon bulge), baseline 5
\addplot[sigblue, thick] {5 + 1.4*exp(-((x-95)^2)/2200)};
% sEMG envelope: muscle burst, baseline 3
\addplot[quotientorange, thick] {3 + 1.2*exp(-((x-78)^2)/700)};
% high-rate accel: sharp onset spike, baseline 1
\addplot[specred, thick] {1 + 1.6*exp(-((x-70)^2)/12)};
\draw[gray, dashed] (axis cs:70,-0.5) -- (axis cs:70,6.5);
\node[gray, anchor=south, font=\footnotesize] at (axis cs:70,6.5) {onset};
\end{axis}
\end{tikzpicture}
\caption{One finger tap across three channels, time-aligned. The high-rate accelerometer (bottom)
spikes first at the contact onset; the sEMG envelope (middle) marks the muscle burst; the FMG
channel (top) rises and falls more slowly as tendon force builds. The fast accelerometer onset is
what lets the live path fire early, while the slower channels confirm which finger and how hard.}
\label{ch40:fig-trace}
\end{figure}

\subsection{The per-finger calibration routine -- the heart of it}

Everything rests on enrollment. Multimodal fusion (high-rate IMU + FMG array + the bio-acoustic
accelerometer + sEMG) gives a richer signature than any single-modality lab paper used --
Figure~\ref{ch40:fig-trace} shows one tap across the channels -- and
per-finger calibration is what turns that signature into a decision. Combining motion with acoustic
contact echoes the IMU-plus-acoustic gesture work (PLOS One, 2020), and wrist-IMU gesture counting
(MDPI Sensors, 2023) shows how much one well-conditioned motion channel already carries; the edge
here is fusing those channels and calibrating per finger rather than leaning on any one of them. The
routine runs on the phone:

\begin{enumerate}
  \item The phone prompts: ``tap your index finger.'' The user does it, say, five times. The wrist
    streams the conditioned channels; the phone slices out each onset.
  \item It \emph{aligns} the repetitions on their onset and \emph{averages} them into one clean
    template $\mathbf{t}_{\text{index}}$ -- a short multi-channel snippet. Averaging suppresses the
    per-trial noise and leaves the shared signature.
  \item It records the spread across repetitions to set a per-finger \emph{threshold}: how close a
    live signature must be to count as a match, and how far from ``idle'' before anything fires.
  \item Repeat for each finger (and for any chord or combined movement you want as its own gesture).
  \item The phone ships the finished templates and thresholds back to the wrist node.
\end{enumerate}

There is no learned model anywhere -- ``analysis'' is filtering, onset alignment, averaging, and
threshold-picking. Calibration is what makes a band that fits \emph{this} arm read \emph{this}
user's fingers, and it is why the method generalizes across people without retraining anything.

\subsection{Live template matching}

On the wrist, live detection is a correlation. With templates $\mathbf{t}_f$ for fingers
$f=1\dots F$ and a live window signature $\mathbf{s}$, pick the finger whose template is closest,
provided the onset clears the idle threshold:
\[
  \hat{f} = \arg\max_{f}\ \frac{\inner{\mathbf{s}}{\mathbf{t}_f}}{\norm{\mathbf{s}}\,\norm{\mathbf{t}_f}},
  \qquad\text{fire only if onset energy} > \theta_{\text{idle}}.
\]
The normalized inner product (cosine similarity) is robust to overall pressure scaling, and the
onset gate keeps resting fingers quiet. For $F\approx 5$--$10$ templates of a few dozen samples
each, this is a handful of dot products -- microseconds on an M4F, trivially inside budget.

\begin{lstlisting}[language=C, caption={Live loop: condition, match against calibrated templates, emit MIDI.}]
void on_block_ready(void) {         // short window of aligned channels
  signature_t s;
  build_signature(&s);              // FMG + sEMG env + accel onset + IMU
  if (onset_energy(&s) > theta_idle) {
    int f = best_template(&s, templ, F);   // argmax cosine similarity
    emit_midi_for_finger(f, pressure(&s)); // Note-On / CC over BLE-MIDI
  }
  block_ready = false;
}

void main_loop(void) {
  for (;;) {
    if (block_ready) on_block_ready();
    enter_low_power_sleep();        // wake on next sensor interrupt
  }
}
\end{lstlisting}

\subsection{Latency budget}

Latency is sample-to-MIDI, and the player feels everything past roughly $20$--$30$~ms. Budget it:

\begin{center}
\small
\begin{tabular}{@{}lr@{}}
\toprule
Stage & Typical \\
\midrule
Sense / fill a short window (8 samples at $200$~Hz) & $\sim 40$~ms \\
Condition + build signature & $< 0.1$~ms \\
Template match (a few dot products) & $< 0.05$~ms \\
BLE-MIDI transmit & $\sim 5$--$10$~ms \\
\midrule
\textbf{Total} & \textbf{$\sim 45$--$50$~ms} \\
\bottomrule
\end{tabular}
\end{center}

The window-fill term dominates, so shorten it: detect on a fast \emph{onset} from the high-rate
accelerometer and sEMG envelope rather than waiting for a full slow-IMU block. Onset detection can
fire within a sample or two of contact, pulling sense-to-decision down to a few milliseconds; the
slower channels then confirm and set pressure. With a short onset window the total lands near
$10$--$15$~ms -- inside the playable range. Choosing the block length is choosing your latency, and
the contact accelerometer is what buys the fast path.

\subsection{MIDI mapping}

\begin{itemize}
  \item \textbf{Note:} each finger maps to a Note-On/Note-Off. A tap is a quick on/off; a hold
    sustains the note while pressure stays above threshold.
  \item \textbf{Velocity:} the onset energy (how hard the tap) maps to MIDI velocity ($0$--$127$).
  \item \textbf{Chords:} a calibrated multi-finger gesture is its own template, fired as several
    Note-Ons together.
  \item \textbf{CC:} continuous FMG pressure or wrist tilt maps to a continuous controller (e.g.
    CC~$74$ for filter cutoff, CC~$1$ for vibrato), so a held note can swell and bend.
\end{itemize}

\subsection{Failure modes}

\begin{itemize}
  \item \textbf{Band shifts on the arm.} Templates assume the same sensor placement as enrollment;
    a big slip detunes them. Mitigation: a quick re-enrollment, and a snug, repeatable band fit.
  \item \textbf{Sweat / electrode drift.} sEMG baseline wanders; lean on FMG and the accelerometer,
    and high-pass the sEMG envelope's baseline.
  \item \textbf{Cross-talk between adjacent fingers.} Neighboring tendons co-activate; the
    multi-channel template and the onset gate separate them better than any single channel.
  \item \textbf{False fires at rest.} Set $\theta_{\text{idle}}$ from the enrollment spread, not by
    guessing.
\end{itemize}

\subsection{The upgrade path}

This build detects which finger moved, how hard, and when -- enough to play. The one thing it
\emph{cannot} do is recover continuous per-joint finger angles from the wrist, because that
information is not in wrist-surface signals (Section~\ref{ch40:compare}). If you later want smooth,
continuous per-finger control -- a finger acting as a fader, not a key -- add finger-worn sensors
(instrumented rings, in the spirit of the ``Learning on the Rings'' work). That is a clean addition,
not a redesign: the rings stream another sample-aligned channel into the same fusion-and-match
pipeline.

Every theme of the chapter is present in this build: hardware-clocked acquisition (no jitter), ring
buffers holding the windows, conditioning in a chosen numeric format, CMSIS-DSP filters and FFTs
well inside budget, a thin BLE-MIDI output, and compute-and-sleep to protect the battery. The same
Fourier and filtering mathematics from the early chapters, now living inside a chip on your arm.

\begin{recoverybox}
If this chapter felt like a gear change, that is because it is one. We left the frictionless
notebook and met deadlines, milliamps, and bits that overflow. The
consolation is that none of the \emph{mathematics} changed -- an FFT is an FFT. What changed is that
you now also account for \emph{when} the work happens and \emph{what it costs}. Carry these three
questions into every embedded design: Is my clock from hardware? Does my work fit the budget? Will
the battery survive? Answer those, and the math takes care of itself.
\end{recoverybox}

\section{Exercises}
\label{ch40:exercises}

\begin{exercise}
A timer-driven sampler uses an AVR with a $16$~MHz clock, Timer1 in CTC mode, and a prescaler of
$64$. What value of \texttt{OCR1A} gives exactly $f_s = 2000$~Hz? (Recall $f_s = f_{\text{clk}} /
(\text{prescaler}\times(\text{OCR1A}+1))$.)
\end{exercise}

\begin{exercise}
Explain in one paragraph why \texttt{delay()}-based sampling can distort a spectrum even when the
average sample rate is correct. Give a concrete example of a code change that would shift the rate
mid-stream.
\end{exercise}

\begin{exercise}
A clock has RMS jitter of $\delta = 1~\mu$s. Using the linearization
$\text{error}\approx 2\pi f A\,\delta$, compute the worst-case amplitude error (as a fraction of
$A$) when sampling a $50$~Hz tone and a $5$~kHz tone. What does the ratio tell you about
high-frequency sampling?
\end{exercise}

\begin{exercise}
You need a ring buffer that holds the last $200$ samples. Why would you allocate $256$ slots
instead of exactly $200$, and what single machine instruction does that choice enable for the
index wrap?
\end{exercise}

\begin{exercise}
In $Q_{1.15}$ format, what integer represents $0.75$? Multiply the Q15 representations of $0.75$
and $0.5$ using the wide-accumulator rule and confirm the result decodes to $0.375$. Show each
step.
\end{exercise}

\begin{exercise}
A Q15 FIR filter has $64$ taps, each up to magnitude $1$, and inputs up to magnitude $1$. Argue why
a $16$-bit accumulator is unsafe and roughly how many guard bits a safe accumulator needs to hold
the worst-case sum without overflow.
\end{exercise}

\begin{exercise}
An $80$~MHz Cortex-M4F runs a $256$-point float FFT in about $30~\mu$s. If you want a new spectrum
every $5$~ms, what fraction of the CPU does the FFT use? Roughly how many such FFTs per frame could
the math alone support, and name two practical limits you would hit first.
\end{exercise}

\begin{exercise}
You sample an IMU at $f_s=200$~Hz, process in $32$-sample blocks, on an $80$~MHz MCU. (a) How many
cycles of budget do you have per block? (b) If your processing costs $40{,}000$ cycles per block,
what fraction of the budget is used? (c) What happens to the budget if you raise $f_s$ to
$2000$~Hz with the same block size?
\end{exercise}

\begin{exercise}
Explain why a Raspberry Pi is a poor choice for an always-on, battery-powered wrist device but an
excellent choice for prototyping the gesture-detection algorithm. Frame your answer in terms of the
two budgets (cycles and joules) from this chapter.
\end{exercise}

\begin{exercise}
Write the three MIDI bytes (in hex) for: (a) a Note-On of middle C (note $60$) at velocity $100$ on
channel $1$ (status nibble $0$); (b) a continuous-controller message setting CC~$74$ to value $64$
on channel~$1$.
\end{exercise}

\begin{exercise}
\textbf{Sensing limits.} A teammate proposes using only the wrist IMU to drive five independent
synth voices, one per finger, by \emph{continuous} finger flexion -- each finger a smooth fader.
Explain precisely why continuous per-joint angles cannot be recovered from a wrist band alone, and
propose what sensor(s) you would add to make continuous per-finger control possible. Separately,
explain why discrete per-finger \emph{detection} (which finger tapped) \emph{is} achievable with the
forearm FMG/sEMG channels plus calibration.
\end{exercise}

\begin{exercise}
\textbf{Calibration.} Describe the per-finger enrollment routine for the capstone: what the phone
prompts, what it records, and how it turns several repetitions of one finger movement into a single
template plus a threshold. Why does averaging aligned repetitions improve the template, and why is
the threshold derived from the spread across repetitions rather than guessed?
\end{exercise}

\begin{exercise}
\textbf{Template match.} Live detection picks $\hat{f}=\arg\max_f
\inner{\mathbf{s}}{\mathbf{t}_f}/(\norm{\mathbf{s}}\norm{\mathbf{t}_f})$ above an onset threshold.
Why use the \emph{normalized} inner product rather than the raw dot product, and what does the onset
gate $\theta_{\text{idle}}$ prevent? For $F=6$ templates of $24$ samples each, estimate the
multiply-accumulate count per match and confirm it is negligible against a $12.8$M-cycle block
budget.
\end{exercise}

\begin{exercise}
\textbf{Design.} Sketch (in prose or pseudocode) the compute-and-sleep loop for the wrist unit, and
estimate the duty cycle if each block takes $50~\mu$s of compute and a block arrives every $160$~ms.
What does that duty cycle imply for average current relative to running the CPU continuously?
\end{exercise}
